\begin{lemma}
\label{editorial-source-lemma-support-closed}
Let $R$ be a ring and let $M$ be an $R$-module.
If $M$ is finite, then $\text{Supp}(M)$ is closed.
More precisely, if $I = \text{Ann}(M)$ is the annihilator of $M$, then
$V(I) = \text{Supp}(M)$.
\end{lemma}

\begin{proof}
We will show that $V(I) = \text{Supp}(M)$.

\medskip\noindent
Suppose $\mathfrak p \in \text{Supp}(M)$. Then $M_{\mathfrak p} \not = 0$.
Choose an element $m \in M$ whose image in $M_\mathfrak p$ is nonzero.
Then the annihilator of $m$ is contained in $\mathfrak p$ by construction
of the localization $M_\mathfrak p$. Hence
a fortiori $I = \text{Ann}(M)$ must be contained in $\mathfrak p$.

\medskip\noindent
If $M=0$, then $I=R$ and both sets are empty.
Otherwise retain generators $x_1,\ldots,x_r$ with $r\geq1$.
Suppose $\mathfrak p\notin\operatorname{Supp}(M)$.
Each $x_i/1$ is zero in $M_{\mathfrak p}$, so
Lemma \ref{algebra-lemma-localization-colimit} gives a common
$f\notin\mathfrak p$ for which all are zero in $M_f$.
For each original generator choose $n_i\geq1$ with
$f^{n_i}x_i=0$, and keep $n=\max\{n_1,\ldots,n_r\}$.
The complete equality
$f^nx_i=f^{n-n_i}(f^{n_i}x_i)=0$ shows that $f^n$ kills
every generator and therefore $f^nM=0$.
Thus $f^n\in I$ while $f^n\notin\mathfrak p$, proving
$\mathfrak p\notin V(I)$ and the required equality.
\end{proof}

\begin{lemma}
\label{editorial-source-lemma-support-finite-annihilator-detection}
Let $M$ be an arbitrary $R$-module. If finitely many
$m_1,\ldots,m_n\in M$ satisfy
$\operatorname{Ann}_R(M)=\bigcap_{i=1}^n\operatorname{Ann}_R(m_i)$,
then $\operatorname{Supp}(M)=V(\operatorname{Ann}_R(M))$.
\end{lemma}

\begin{proof}
Retain the original submodule $N=\sum_{i=1}^n Rm_i\subseteq M$.
An element kills $N$ exactly when it kills each $m_i$,
so $\operatorname{Ann}_R(N)=\operatorname{Ann}_R(M)$ by
the stated hypothesis. Since $N$ is finite, the preceding
lemma gives $\operatorname{Supp}(N)=V(\operatorname{Ann}_R(M))$.
Localizing the actual injection $N\hookrightarrow M$
is injective, so $\operatorname{Supp}(N)\subseteq\operatorname{Supp}(M)$.
For every $M$, the first half of the preceding proof
gives $\operatorname{Supp}(M)\subseteq V(\operatorname{Ann}_R(M))$.
The inclusions prove equality without replacing the
original $M$ by $N$. If the list is empty its intersection
of annihilators is $R$, hence $M=0$, which also satisfies
the conclusion. For example an arbitrary nonempty free
module is covered by this hypothesis: one original basis
vector already has annihilator zero, although the basis
may be infinite.
\end{proof}